\title{DeepSeekMath: Pushing the Limits of Mathematical Reasoning in Open Language Models}

\begin{abstract}
The intricate and highly structured nature of mathematical reasoning presents considerable difficulties for language models. Here, we present DeepSeekMath~7B, developed by continuing the pre-training of DeepSeek-Coder-Base-v1.5 7B using a dataset comprising 120B math-related tokens derived from Common Crawl, supplemented with both natural language and code corpora. DeepSeekMath~7B attains a remarkable 51.7\% on the MATH benchmark, designed for competition-level mathematical reasoning, without the need for external toolkits or ensemble methods, thus narrowing the gap with models such as Gemini-Ultra and GPT-4. Employing a self-consistency approach over 64 samples, DeepSeekMath~7B achieves 60.9\% on the same benchmark. This elevated mathematical reasoning performance can be ascribed to two principal innovations: first, leveraging large-scale web-sourced data via a carefully crafted data curation pipeline; second, the introduction of Group Relative Policy Optimization (GRPO)—an enhanced version of Proximal Policy Optimization (PPO)—which not only bolsters mathematical problem-solving skills but also improves PPO’s memory efficiency.
\end{abstract}

\section{Introduction}

Large language models (LLMs) have fundamentally transformed methods for mathematical reasoning within artificial intelligence, driving notable progress across quantitative reasoning \citep{MATH} and geometric reasoning benchmarks \citep{trinh2024solving}. In addition, these systems serve as valuable tools to aid humans in resolving sophisticated mathematical challenges \citep{tao}. Nevertheless, leading solutions such as GPT-4 \citep{gpt4} and Gemini-Ultra \citep{gemini} remain inaccessible to the public, resulting in a substantial performance disparity, as currently available open-source models lag considerably behind these proprietary counterparts.

This work presents DeepSeekMath, a specialized language model that markedly surpasses existing open-source models in mathematical reasoning and approaches the benchmark performance of GPT-4 in academic evaluations. Central to this advancement is the development of the DeepSeekMath Corpus, a substantial and high-quality pre-training resource containing 120B math-focused tokens. The corpus is derived from Common Crawl (CC) by leveraging a fastText-based classifier \citep{joulin2016fasttext}. For its initial training, the classifier utilizes examples from OpenWebMath \citep{openwebmath} as positives, while negatives consist of a varied sample of unrelated web content. Using this model, additional positive samples are identified within the CC and subsequently curated via human annotation, after which the classifier is further refined using this improved set. Results from our experiments demonstrate that the corpus achieves strong quality, as evidenced by the performance of our baseline DeepSeekMath-Base 7B model, which obtains 64.2\% accuracy on GSM8K \citep{gsm8k} and 36.2\% on the advanced MATH dataset \citep{MATH}, surpassing the results of Minerva 540B \citep{minerva}. Additionally, because the DeepSeekMath Corpus encompasses multiple languages, we observe notable gains on Chinese mathematics benchmarks \citep{wei2023cmath,agieval}. We suggest that our methodologies for mathematical data curation can serve as a valuable foundation for future research, with ample prospects for further enhancements.

DeepSeekMath-Base is built upon the DeepSeek-Coder-Base-v1.5 7B \citep{deepseek-coder}, as our findings indicate that initiating model development from a code-focused pre-trained model yields better results than starting from a general LLM. Additionally, we note that further math-specific training enhances not only the model’s mathematical competence but also its general reasoning abilities, as demonstrated by performance improvements on MMLU \citep{mmlu} and BBH benchmarks \citep{bbh}.

Following pre-training, we perform mathematical instruction tuning on DeepSeekMath-Base utilizing datasets tailored to chain-of-thought \citep{cot}, program-of-thought \citep{pot,pal}, and tool-integrated reasoning \citep{tora}. The tuned model, DeepSeekMath-Instruct 7B, surpasses all other models at the 7B parameter scale and achieves performance on par with leading 70B open-source instruction-tuned systems.

Moreover, we present Group Relative Policy Optimization (GRPO), a modified reinforcement learning (RL) approach derived from Proximal Policy Optimization (PPO) \citep{schulman2017proximal}. Unlike traditional methods, GRPO eliminates the need for a critic network by leveraging group-based score baselines, substantially decreasing the computational burden required for training. Using only a selection of English instruction tuning data, GRPO secures significant advancements over the already strong DeepSeekMath-Instruct model, with gains in both in-domain (GSM8K: 82.9\% $\rightarrow$ 88.2\%, MATH: 46.8\% $\rightarrow$ 51.7\%) and out-of-domain tasks (e.g., CMATH: 84.6\% $\rightarrow$ 88.8\%) throughout reinforcement learning. We further propose a unified framework to interpret various techniques, encompassing Rejection Sampling Fine-Tuning (RFT) \citep{yuan2023scaling}, Direct Preference Optimization (DPO) \citep{dpo}, PPO, and GRPO. Within this conceptualization, these strategies can all be viewed as either straightforward or streamlined RL variants. Comprehensive experiments are carried out—comparing online versus offline training, outcome versus process-based supervision, and single-turn versus iterative RL—to thoroughly dissect critical components of this framework. Finally, we elucidate how RL enhances instruction-tuned model performance and outline prospective directions for advancing RL under this unified perspective.

\subsection{Contributions}
This work presents advances in large-scale math pre-training and provides a systematic investigation into reinforcement learning.

\textbf{Math Pre-Training at Scale}

\begin{itemize}[topsep=0pt]
    \item We demonstrate that the open-source Common Crawl dataset serves as a rich resource for mathematics. By employing a carefully engineered pipeline for data curation, we assemble the DeepSeekMath Corpus—a refined, high-quality dataset comprising 120B tokens extracted from web pages identified as containing mathematical material. This dataset is nearly 7 times larger than the collection used by Minerva \citep{minerva} and exceeds the size of the recently introduced OpenWebMath \citep{openwebmath} by a factor of 9.

    \item The DeepSeekMath-Base 7B model, our pre-trained foundation model, attains performance on par with the much larger Minerva 540B \citep{minerva}. This suggests that, in mathematical reasoning tasks, model scale is not the sole determinant of ability—a compact model, when trained on sufficiently curated data, can yield robust results.

    \item We report experimental insights concerning math training. Pre-training models on code data prior to math-focused training enhances their proficiency in solving mathematical tasks, with and without external tools. This observation provides partial resolution to the enduring question: \textit{does code training benefit reasoning skills?} Our results support its positive impact, particularly in mathematical domains.

    \item Despite the prevalent practice of leveraging arXiv papers for training, our evaluations indicate that this approach does not deliver significant performance gains across the mathematical benchmarks considered in this study.
\end{itemize}

\textbf{Exploration and Analysis of Reinforcement Learning}

\begin{itemize}[topsep=0pt]
    \item We present Group Relative Policy Optimization (GRPO), a novel reinforcement learning method that operates without relying on a critic model. By estimating the baseline via aggregated group scores, GRPO substantially decreases the computational demands of training when compared to Proximal Policy Optimization (PPO).
    \item Our empirical results reveal that GRPO notably boosts the capabilities of our instruction-tuned DeepSeekMath-Instruct model, utilizing only the instruction-tuning dataset. Additionally, GRPO fosters improved generalization, as reflected in superior out-of-domain performance during reinforcement learning.
    \item We develop a cohesive framework that connects a range of techniques, including RFT, DPO, PPO, and GRPO. To systematically analyze this framework, we conduct a comprehensive suite of experiments, such as contrasting online and offline learning, evaluating the roles of outcome versus process supervision, and comparing single-turn with iterative reinforcement learning, among others.
    \item Building on the established unified framework, we probe the mechanisms driving the success of reinforcement learning, and propose several promising directions for advancing the reinforcement learning capabilities of large language models (LLMs).
\end{itemize}

\subsection{Summary of Evaluations and Metrics}

\begin{itemize}[topsep=0pt]
    \item \textbf{English and Chinese Mathematical Reasoning}: We carry out thorough evaluations on both English and Chinese datasets, spanning mathematics problems from elementary to collegiate levels. The English evaluations utilize datasets such as GSM8K \citep{gsm8k}, MATH \citep{MATH}, SAT \citep{llemma}, OCW Courses \citep{minerva}, and MMLU-STEM \citep{mmlu}. Chinese tasks are benchmarked with MGSM-zh \citep{mgsm}, CMATH \citep{wei2023cmath}, Gaokao-MathCloze \citep{agieval}, and Gaokao-MathQA \citep{agieval}. Our evaluation covers both the ability of models to construct detailed, stand-alone textual solutions and the capacity to address problems computationally using Python.

    DeepSeekMath-Base achieves results comparable to the proprietary Minerva 540B \citep{minerva} on English math tasks and surpasses all publicly available base models (such as Mistral 7B \citep{mistral} and Llemma-34B \citep{llemma}), with or without mathematical pre-training, often with a substantial lead. For Chinese benchmarks, DeepSeekMath-Base further distinguishes itself, potentially due to our inclusion of high-quality non-English pre-training data, diverging from previous efforts \citep{minerva,llemma} that focus solely on English datasets. Upon applying mathematical instruction tuning and reinforcement learning, DeepSeekMath-Instruct and DeepSeekMath-RL attain robust performance, achieving, for the first time in open-source models, accuracy rates exceeding 50\% on the challenging, competition-level MATH dataset.

\item \textbf{Formal Mathematics}: To assess the capabilities of DeepSeekMath-Base, we apply it to the informal-to-formal theorem proving task described in \citep{dsp_proof}, utilizing miniF2F \citep{minif2f} as the evaluation benchmark and Isabelle \citep{isabelle} as the proof assistant. DeepSeekMath-Base achieves notable results in few-shot autoformalization.

\item \textbf{Natural Language Understanding, Reasoning, and Code}: In order to gain a broad understanding of the model's proficiency in general comprehension, reasoning, and coding, we benchmark DeepSeekMath-Base on Massive Multitask Language Understanding (MMLU) \citep{mmlu}, which spans 57 multiple-choice domains, BIG-Bench Hard (BBH) \citep{bbh}, which contains 23 problems emphasizing multi-step reasoning, and evaluate its code generation abilities using HumanEval \citep{codex} and MBPP \citep{mbpp}. The results indicate that mathematical pre-training enhances both reasoning and general language understanding capabilities.

\section{Math Pre-Training}

\subsection{Data Collection and Decontamination} 

Here, we describe the methodology for assembling the DeepSeekMath Corpus leveraging data from Common Crawl. Illustrated in Figure \ref{fig:data_collect}, our approach utilizes an iterative framework designed to efficiently expand a large-scale mathematical text collection by beginning with a high-quality, smaller seed dataset of mathematical resources. Notably, this workflow can be extended to other specialized domains, including computer code.

Our procedure starts by designating OpenWebMath \citep{openwebmath}—a repository of rigorously curated mathematical web content—as the initial seed. Leveraging this set, we train a fastText classifier \citep{joulin2016fasttext} to identify and extract additional mathematical webpages with characteristics similar to those in OpenWebMath. We draw 500,000 samples from the seed as positive training examples and randomly select 500,000 pages from Common Crawl to serve as negatives. For training, we utilize an open-source fastText implementation,\footnote{\url{https://fasttext.cc}} setting the embedding dimension to 256, learning rate to 0.1, word n-gram maximum length to 3, minimum word frequency to 3, and running for 3 training epochs. To mitigate the scale of the raw Common Crawl corpus, we employ URL-based deduplication and near-duplicate filtering, resulting in a reduced set of 40B unique HTML pages. Utilizing the trained fastText model, we identify and extract pages likely to be mathematical in content from the deduplicated dataset. Candidate pages are ranked based on the classifier’s output scores, retaining only those with the highest predicted relevance. The appropriate amount of data to preserve is informed by pre-training performance over datasets sized at 40B, 80B, 120B, and 160B tokens, and in this initial phase, we select the top 40B tokens for retention.

Following the initial round of data collection, a substantial number of mathematical web pages remain uncaptured. This limitation primarily stems from the fastText model, which has been trained on a set of positive samples that lacks adequate variety. To enhance the fastText model’s effectiveness, we identify and incorporate additional mathematical web domains to further enrich the initial seed corpus. Our methodology begins with segmenting the entire Common Crawl dataset into non-overlapping domains, where each domain comprises web pages sharing an identical base URL. For each domain, we assess the fraction of pages gathered during the first round of collection. Any domain with more than 10\% of its pages already collected is designated as math-oriented (for instance, \url{mathoverflow.net}).

Next, we manually review URLs within these recognized domains to annotate those associated with mathematical topics (e.g., \url{mathoverflow.net/questions}). The web pages tied to these newly annotated URLs, if not previously collected, are added to our seed corpus. This strategy enables us to compile a broader set of positive examples and, subsequently, to train a more robust fastText model with enhanced recall for mathematical content in subsequent data collection rounds. Across four iterations, we ultimately collect 35.5M math-related web pages, amassing a total of 120B tokens. During the fourth round, we observe that roughly 98\% of the content had already been acquired by the third round, leading us to conclude the data collection process.

To prevent test set leakage, we apply filtering procedures similar to those detailed in \cite{deepseek-coder}: we remove all web pages that contain questions or solutions from English math benchmarks like GSM8K \citep{gsm8k} and MATH \citep{MATH}, as well as from Chinese benchmarks such as CMATH \citep{wei2023cmath} and AGIEval \citep{agieval}. The filtering rule is that any segment of text featuring an exact match to any 10-gram phrase from these benchmark datasets is excluded from our training collection. For benchmark excerpts between 3 and 9 grams in length, we also apply exact matching to identify and eliminate any contaminated web pages.

\subsection{Validating the Quality of the DeepSeekMath~Corpus}

We conduct pre-training experiments to assess the DeepSeekMath~Corpus in relation to the most recent publicly available math datasets:

\begin{itemize}[topsep=0pt]
    \item \textbf{MathPile} \citep{mathpile}: a composite dataset (8.9B tokens) drawn from multiple sources such as textbooks, Wikipedia, ProofWiki, CommonCrawl, StackExchange, and arXiv, with more than 85\% of content sourced from arXiv;
    \item \textbf{OpenWebMath} \citep{openwebmath}: a dataset filtered from CommonCrawl for mathematical material, amounting to 13.6B tokens;
    \item \textbf{Proof-Pile-2} \citep{llemma}: a collection combining OpenWebMath, AlgebraicStack (comprising 10.3B tokens of mathematical code), and arXiv papers (28.0B tokens). For experiments involving Proof-Pile-2, we adopt the arXiv:Web:Code proportion of 2:4:1 as proposed by \cite{llemma}.
\end{itemize}

\subsubsection{Training Setting}
\label{sec:quality-policy}
We perform specialized math-oriented training on a general-purpose language model with 1.3 billion parameters, which is architecturally consistent with the DeepSeek LLMs \citep{deepseek-llm}, and referred to as DeepSeek-LLM 1.3B in this study. Each mathematical dataset is used to independently train a model for 150 billion tokens. All training is executed using the resource-efficient HAI-LLM framework \citep{haillm}. In line with the established DeepSeek LLMs methodology, we utilize the AdamW optimizer \citep{adamW} configured with $\beta_1=0.9$, $\beta_2=0.95$, and a $\mathrm{weight\_decay}$ of 0.1. The learning rate schedule is multi-phased: after an initial 2,000-step warmup period, the rate peaks, drops to 31.6\% at 80\% completion of training, and further falls to 10.0\% of the maximum by 90\% of training. The peak learning rate is set at 5.3e-4, and we train with a global batch size of 4 million tokens and a context window of 4,000 tokens.

\subsubsection{Evaluation Results}

\textbf{The DeepSeekMath~Corpus distinguishes itself by its superior quality, extensive multilingual coverage, and unmatched scale.}

\begin{itemize}[topsep=0pt]
    \item \textbf{High-quality}:
    To assess model quality, we examine downstream results across 8 mathematical evaluation sets employing few-shot chain-of-thought prompting \cite{cot}. As detailed in Table \ref{tab:corpora_comparison}, the model trained on DeepSeekMath~Corpus substantially surpasses alternatives. Figure \ref{fig:corpora_comparisons} highlights that, after training on just 50B tokens (equivalent to a single Proof-Pile-2 epoch), DeepSeekMath-based training yields superior outcomes, suggesting that DeepSeekMath~Corpus exhibits higher average quality compared to Proof-Pile-2.
    \item \textbf{Multilingual}:
    The DeepSeekMath~Corpus contains mathematical material in several languages, with English and Chinese being most prevalent. Table \ref{tab:corpora_comparison} indicates that models trained on this corpus achieve elevated mathematical reasoning abilities in both English and Chinese. By comparison, other available mathematical corpora—primarily composed of English text—confer little benefit and may in fact compromise Chinese mathematical reasoning capabilities.
    \item \textbf{Large-scale}:
    In terms of data volume, DeepSeekMath~Corpus is several times larger than any existing mathematical dataset. Figure \ref{fig:corpora_comparisons} demonstrates that, when DeepSeek-LLM 1.3B is trained on DeepSeekMath~Corpus, the learning curve is steeper and improvement persists longer. In contrast, smaller baseline corpora are exhausted after only a few training repetitions, leading model performance to plateau rapidly.
\end{itemize}

\subsection{Training and Evaluating DeepSeekMath-Base 7B}

This section presents DeepSeekMath-Base 7B, a foundational model engineered for advanced reasoning, with particular proficiency in mathematical contexts. The initialization utilizes DeepSeek-Coder-Base-v1.5 7B \citep{deepseek-coder}, and training is carried out for 500B tokens. The data composition comprises 56\% from the DeepSeekMath Corpus, 4\% sourced from AlgebraicStack, 10\% derived from arXiv, 20\% constituted by Github code repositories, and the remaining 10\% originating from Common Crawl, with natural language text in both English and Chinese. Training generally adheres to the protocols described in Section \ref{sec:quality-policy}, apart from setting the maximum learning rate at 4.2e-4 and employing a batch size of 10M tokens.

A thorough evaluation of DeepSeekMath-Base 7B is undertaken to analyze its mathematical reasoning, particularly its capacity for generating independent, self-contained solutions, utilizing computational tools for problem-solving, and executing formal theorem proving. In addition, a broader examination of the model's capabilities is included, covering natural language understanding, general reasoning aptitude, and programming proficiency.

\paragraph{Mathematical Problem Solving with Step-by-Step Reasoning}

DeepSeekMath-Base is assessed on its aptitude for mathematical problem-solving using few-shot chain-of-thought prompting \citep{cot}, spanning eight benchmarks in both English and Chinese. These benchmarks address quantitative reasoning tasks (such as GSM8K \citep{gsm8k}, MATH \citep{MATH}, and CMATH \citep{wei2023cmath}), as well as multiple-choice question sets (for example, MMLU-STEM \citep{mmlu} and Gaokao-MathQA \citep{agieval}). The evaluation encompasses a diverse array of mathematical domains, ranging from elementary through to collegiate-level challenges.

As illustrated in Table \ref{tab:base_math_cot}, DeepSeekMath-Base 7B consistently achieves superior results across all eight assessed benchmarks when compared with open-source base models—including the popular generalist Mistral 7B \citep{mistral} and the newly released, math-focused Llemma 34B \citep{llemma}, both of which received specialized mathematical training (Llemma on Proof-Pile-2 \citep{llemma}). In particular, DeepSeekMath-Base demonstrates a margin exceeding 10\% absolute improvement over previous open-source base models on the challenging, competition-grade MATH dataset. Furthermore, it outperforms Minerva 540B \citep{minerva}, a proprietary model based on PaLM \citep{palm} that is more than 77 times its size and further fine-tuned on mathematical content.

\paragraph{Mathematical Problem Solving with Tool Use} For program-assisted mathematical reasoning, we conduct evaluations on GSM8K and MATH using few-shot program-of-thought prompting \citep{pot,pal}. Here, models approach each problem by constructing a Python script, with access to libraries such as \textit{math} and \textit{sympy} for complex calculations. The produced program is executed and its output is taken as the solution. Table \ref{tab:base_math_pot_proof} indicates that DeepSeekMath-Base 7B achieves performance that surpasses the previously leading Llemma 34B model.

\paragraph{Formal Mathematics}

Automation in formal proof construction enhances both the rigor and efficiency of mathematical reasoning, and has seen growing interest in recent times. We assess DeepSeekMath-Base 7B on the informal-to-formal proof translation task defined in \citep{dsp_proof}: the goal is to produce a formalized proof using an informal problem statement, its formal equivalent, and an accompanying informal proof. Experiments are conducted on miniF2F \citep{minif2f}, a benchmark targeting Olympiad-level formal mathematics. For each instance, the model generates an Isabelle proof with few-shot prompts. Following the methodology of \cite{dsp_proof}, the model drafts a proof sketch, with the automated prover Sledgehammer \citep{sledgehammer} subsequently resolving any remaining gaps. Table \ref{tab:base_math_pot_proof} shows that DeepSeekMath-Base 7B performs robustly on proof autoformalization.

\paragraph{Natural Language Understanding, Reasoning, and Code}

To gauge performance in related domains, we evaluate on MMLU \citep{mmlu} for natural language understanding, BBH \citep{bbh} for reasoning, and HumanEval \citep{codex} along with MBPP \citep{mbpp} for coding proficiency. The results in Table \ref{tab:understanding_reasoning_code} reveal that DeepSeekMath-Base 7B makes substantial advances on both MMLU and BBH relative to its predecessor DeepSeek-Coder-Base-v1.5 \citep{deepseek-coder}, underscoring the beneficial effects of mathematical pretraining on comprehension and reasoning. The continued incorporation of code tokens during training also allows DeepSeekMath-Base 7B to match the coding performance of DeepSeek-Coder-Base-v1.5 across both benchmarks. Taken together, DeepSeekMath-Base 7B offers notable improvements over the generalist Mistral 7B \citep{mistral} on all three tasks involving reasoning and code generation.

\section{Supervised Fine-Tuning}

\subsection{SFT Data Curation}

We assemble an instruction-tuning dataset for mathematics that encompasses both English and Chinese problems from a broad array of mathematical domains and levels of difficulty. Each problem is accompanied by its corresponding solution, utilizing formats such as chain-of-thought (CoT) \citep{cot}, program-of-thought (PoT) \citep{pot,pal}, or tool-integrated reasoning \citep{tora}. In total, this yields 776K training instances.

\begin{itemize}[topsep=0pt]
    \item \textbf{English mathematical datasets}: For the English component, we supplement GSM8K and MATH tasks with tool-integrated answers, select items from MathInstruct \citep{MathInstruct}, and include problems from the Lila-OOD \citep{lila} training corpus where solutions are provided in either CoT or PoT style. This curated collection spans numerous mathematical disciplines such as algebra, probability, number theory, calculus, and geometry.
    \item \textbf{Chinese mathematical datasets}: For the Chinese set, we aggregate K-12 mathematics problems across 76 different subfields (e.g., linear equations), and supply annotated solutions using both CoT and tool-integrated reasoning paradigms.
\end{itemize}

\subsection{Training and Evaluating DeepSeekMath-Instruct 7B}

Here, we detail the fine-tuning of DeepSeekMath-Instruct 7B, which is based on DeepSeekMath-Base and further trained on mathematical instruction data. Training examples are concatenated randomly, up to a limit of 4K tokens per sequence. The model is optimized for 500 steps, utilizing a batch size of 256 and a fixed learning rate of 5e-5.

For model assessment, we gauge mathematical problem-solving capabilities both with and without tool integration on four different quantitative reasoning benchmarks in English and Chinese. The evaluation pits our model against top contemporaneous models:

\begin{itemize}[topsep=0pt]
    \item \textbf{Closed-source models} under comparison are: (1) the GPT series, specifically GPT-4 \citep{gpt4} and GPT-4 Code Interpreter~\footnote{\url{https://openai.com/blog/chatgpt-plugins\#\#code-interpreter}}, (2) Gemini Ultra and Pro \citep{gemini}, (3) Inflection-2 \citep{inflection-2}, (4) Grok-1~\footnote{\url{https://x.ai/model-card}}, in addition to leading offerings from Chinese developers such as (5) Baichuan-3~\footnote{\url{https://www.baichuan-ai.com}}, (6) the recent GLM-4~\footnote{\url{https://open.bigmodel.cn/dev/api\#glm-4}} from the GLM lineage \citep{glm}.
\end{itemize}

These models are designed for broad application scenarios and typically undergo multiple rounds of alignment refinement.

\item \textbf{Open-source models} span both general-purpose and mathematically enhanced variants: the general models include (1) DeepSeek-LLM-Chat 67B \citep{deepseek-llm}, (2) Qwen 72B \citep{qwen}, (3) SeaLLM-v2 7B \citep{seallm}, and (4) ChatGLM3 6B \citep{chatglm3}. Mathematical enhancements are present in models such as (5) InternLM2-Math 20B~\footnote{\url{https://github.com/InternLM/InternLM-Math}}, which is derived from InternLM2 with targeted mathematical training and subsequent instruction tuning, (6) Math-Shepherd-Mistral 7B, developed by applying PPO \citep{schulman2017proximal} to Mistral 7B \citep{mistral} with a reward model supervising the process, (7) the WizardMath suite \citep{wizardmath}, which advances mathematical inference for Mistral 7B and Llama-2 70B \citep{llama2} through an “evolve-instruct” method (leveraging instruction tuning using AI-generated prompts) and PPO training primarily on GSM8K and MATH datasets, (8) MetaMath 70B \citep{metamath}, created by fine-tuning Llama-2 70B with extended GSM8K and MATH data, (9) ToRA 34B \cite{tora}, where CodeLlama 34B is further trained for mathematical reasoning that incorporates tool usage, and (10) MAmmoTH 70B \citep{MathInstruct}, a version of Llama-2 70B fine-tuned with MathInstruct data through instruction tuning.
\end{itemize}

Referencing Table \ref{tab:sft_rl_math}, under conditions where use of external tools is not permitted, DeepSeekMath-Instruct 7B exhibits exceptional capabilities in sequential reasoning. In particular, for the MATH dataset at the competition level, this model outperforms all open-source competitors and nearly all proprietary offerings (such as Inflection-2 and Gemini Pro) by a margin of at least 9\% in absolute terms, even outperforming models with much larger parameter counts (like Qwen 72B) and those utilizing math-specific reinforcement learning strategies (like WizardMath-v1.1 7B). While DeepSeekMath-Instruct attains results comparable to the Chinese proprietary systems GLM-4 and Baichuan-3 on MATH, it does not reach the performance levels of GPT-4 and Gemini Ultra.

In evaluation scenarios permitting the combination of natural language reasoning with programmatic tools, DeepSeekMath-Instruct 7B achieves nearly 60\% accuracy on the MATH benchmark, exceeding the results of all other available open-source models. Across other assessment datasets, it rivals DeepSeek-LLM-Chat 67B, the previously leading model in this domain, despite being an order of magnitude smaller.

\section{Reinforcement Learning}

\subsection{Group Relative Policy Optimization}
Reinforcement learning (RL) has demonstrated notable improvements in LLMs’ mathematical reasoning skills after the completion of the Supervised Fine-Tuning (SFT) phase \citep{wang2023math,wizardmath}. This section introduces Group Relative Policy Optimization (GRPO), a novel and efficient RL technique.

\subsubsection{From PPO to GRPO}
Proximal Policy Optimization (PPO) \citep{schulman2017proximal}, a prominent actor-critic RL framework, is frequently adopted in the fine-tuning of LLMs via reinforcement learning \citep{ouyang2022training}. The PPO objective guides optimization by maximizing the following surrogate function:

\begin{equation}
\footnotesize
    \mathcal{J}_{PPO}(\theta) = \mathbb{E}{[q \sim P(Q), o \sim \pi_{\theta_{old}}(O|q)]} \frac{1}{|o|} \sum_{t=1}^{|o|} \min \left[ \frac{\pi_\theta(o_{t} | q, o_{<t})}{\pi_{\theta_{old}}(o_{t} | q, o_{<t})} A_{t}, \text{clip} \left( \frac{\pi_\theta(o_{t} | q, o_{<t})}{\pi_{\theta_{old}}(o_{t} | q, o_{<t})}, 1 - \vare

Here, $\pi_{\theta}$ denotes the current policy model, while $\pi_{\theta_{old}}$ represents the previous policy. The variables $q$ and $o$ correspond to questions and outputs, which are sampled from both the question dataset and the preceding policy $\pi_{\theta_{old}}$. The parameter $\varepsilon$ acts as a clipping hyper-parameter within PPO, aiming to enhance the stability of the training process. The advantage term, $A_t$, is estimated through Generalized Advantage Estimation (GAE) \citep{gae}, utilizing future rewards $\{r_{\ge t}\}$ in conjunction with a learned value function $V_{\psi}$. Consequently, PPO requires joint training of the value function alongside the policy model. To prevent excessive optimization towards the reward model, it is common practice to impose a per-token KL-divergence penalty with respect to a reference model, added to the reward at each token position, as outlined in \citep{ouyang2022training}:

\begin{equation}
    r_{t} = r_\varphi(q, o_{\le t}) - \beta \log\frac{\pi_{\theta}(o_{t}|q, o_{<t})}{\pi_{ref}(o_{t}|q, o_{<t})},
\label{eq:PPO-reward}
\end{equation}

In this expression, $r_\varphi$ signifies the reward model, $\pi_{ref}$ serves as the reference model (typically initialized as the SFT model), and $\beta$ specifies the weighting for the KL penalty term.

Training a value function in PPO typically involves a model of size comparable to the main policy, which significantly increases memory usage and computation time. Moreover, the value function is leveraged as a baseline for reducing variance in advantage estimates during RL optimization. This becomes especially challenging for LLMs, where only the terminal token is often assigned a reward by the reward model, making it nontrivial to accurately train a value function for every token position. To circumvent these issues, we introduce Group Relative Policy Optimization (GRPO), illustrated in Figure \ref{fig:grpo}, which eliminates the necessity for learning an explicit value function, as is required in PPO. Instead, GRPO adopts the mean reward from a set of outputs—each produced for the same prompt—as the baseline. Concretely, for each question $q$, a batch of outputs $\{o_1, o_2, \cdots, o_G\}$ is generated from $\pi_{\theta_{old}}$. The policy update then maximizes the following objective:

\begin{equation}
\footnotesize
\begin{split}
    \mathcal{J}_{GRPO}(\theta) &= \mathbb{E}{[q \sim P(Q), \{o_i\}_{i=1}^G \sim \pi_{\theta_{old}}(O|q)]}  \\
    & \frac{1}{G}\sum_{i=1}^G\frac{1}{|o_i|} \sum_{t=1}^{|o_i|} \left\{ \min \left[ \frac{\pi_\theta(o_{i,t} | q, o_{i,<t})}{\pi_{\theta_{old}}(o_{i,t} | q, o_{i,<t})} \hat{A}_{i,t}, \text{clip} \left( \frac{\pi_\theta(o_{i,t} | q, o_{i,<t})}{\pi_{\theta_{old}}(o_{i,t} | q, o_{i,<t})}, 1 - \varepsilon, 1 + \varepsilon \right)  \hat{A}_{i,t} \right] - \beta \mathbb{D}_{KL}\left[\pi_{\theta} || \pi_{ref}\right]\right\} ,
\end{split}
\label{eq:GRPO-obj}
\end{equation}

Here, the hyper-parameters $\varepsilon$ and $\beta$ function analogously as before, and $\hat{A}_{i,t}$ denotes the advantage derived solely from relative rewards within each group—details of which are elaborated in the subsequent sections. This group-relative calculation of advantages closely matches the comparison-based construction of reward models, since these are often trained on datasets comparing candidate outputs for the same question. Furthermore, rather than incorporating the KL penalty directly within the reward as in PPO, GRPO imposes the KL-regularization term by explicitly including the KL divergence between the new policy and reference policy in the loss objective, thereby streamlining the computation of $\hat{A}_{i,t}$. Distinct from the KL-penalty formulation used in (\ref{eq:P

\begin{algorithm}[t]
  \small
  \caption{Iterative Group Relative Policy Optimization}
  \textbf{Input} initial policy model $\pi_{\theta_{\text{init}}}$; reward models $r_\varphi$; task prompts $\mathcal{D}$; 
  hyperparameters $\varepsilon$, $\beta$, $\mu$
  \begin{algorithmic}[1]
    \State Initialize the policy model: $\pi_\theta \leftarrow \pi_{\theta_{\text{init}}}$
    \For{iteration = 1, \dots, I}
       \State Set the reference model: $\pi_{ref} \leftarrow \pi_{\theta}$
      \For{step = 1, \dots, M}
      \State Draw a mini-batch $\mathcal{D}_b$ from $\mathcal{D}$
      \State Assign the old policy model: $\pi_{\theta_{old}} \leftarrow \pi_{\theta}$ 
      \State For each prompt $q \in \mathcal{D}_b$, generate $G$ candidate responses $\{o_i\}_{i=1}^G$ using $\pi_{\theta_{old}} (\cdot \mid q) $
      \State Compute corresponding rewards $\{r_i\}_{i=1}^{G}$ for every output $o_i$ via the reward model $r_\varphi$ 
      \State Derive group-relative advantage $\hat{A}_{i,t}$ for the $t$-th token in each output $o_i$
      \For{GRPO iteration = 1, \dots, $\mu$}
        \State Update $\pi_{\theta}$ by maximizing the GRPO criterion (see Equation \ref{eq:GC-GRPO})
      \EndFor
    \EndFor 
    \State Continue updating $r_\varphi$ through ongoing training utilizing a replay buffer.
    \EndFor 
  \end{algorithmic}
  \textbf{Output} $\pi_\theta$
  \label{alg:iter-grpo}
\end{algorithm}

\subsubsection{Outcome Supervision RL with GRPO} More precisely, for a given prompt $q$, one samples a set of responses $\{o_1, o_2, \cdots, o_G\}$ from the prior policy $\pi_{\theta_{old}}$. Each of these outputs receives a score from a reward model, producing the reward vector $\mathbf{r} = \{r_1, r_2, \cdots, r_G\}$. The rewards are standardized by subtracting their mean and dividing by the standard deviation across the group. Outcome supervision attaches this normalized reward to every token in each output $o_i$, so that all advantages are set to the standardized value: $\hat{A}_{i, t} = \widetilde{r}_i = \frac{r_i- {\rm mean}(\mathbf{r})}{{\rm std}(\mathbf{r})}$. The policy is then improved by maximizing the objective presented in equation (\ref{eq:GRPO-obj}).

\subsubsection{Process Supervision RL with GRPO} Since outcome supervision only offers a terminal reward for each output, it can lack granularity and efficiency, especially in settings involving complex mathematical reasoning. Building on \cite{wang2023math}, process supervision is utilized, granting feedback at each intermediate reasoning step. Specifically, for a prompt $q$ with $G$ generated outputs $\{o_1, o_2, \cdots, o_G\}$, a process reward model evaluates each step, providing rewards as $\mathbf{R} = \{ \{r_1^{index(1)},\cdots,r_1^{index(K_1)}\}, \cdots,  \{r_G^{index(1)},\cdots,r_G^{index(K_G)}\} \}$, where $index(j)$ indicates the position of the $j$-th step's final token, and $K_i$ counts steps in output $o_i$. These step-level rewards are likewise standardized: $\widetilde{r}_i^{index(j)} = \frac{r_i^{index(j)} - {\rm mean}(\mathbf{R})}{{\rm std}(\mathbf{R})}$. 
Process supervision then computes the token-level advantage as the sum of standardized rewards of subsequent steps: $\hat{A}_{i, t} = \sum_{index(j) \ge t} \widetilde{r}_i^{index(j)}$. Policy optimization proceeds by maximizing the same objective as specified in equation (\ref{eq:GRPO-obj}).

\subsubsection{Iterative RL with GRPO} During reinforcement learning, the initial reward model may gradually lose its effectiveness as a supervisor for the evolving policy model. To address this, we implement an iterative RL scheme utilizing GRPO. As illustrated in Algorithm \ref{alg:iter-grpo}, the iterative GRPO procedure repeatedly constructs new reward model training sets by sampling from the latest policy model. The reward model is then incrementally updated, leveraging a replay buffer containing 10\% of earlier training samples to retain previous learning. After each iteration, we designate the current policy model as the new reference and proceed to update the policy model with feedback from the newly trained reward model.

\subsection{Training and Evaluating DeepSeekMath-RL}

Our reinforcement learning experiments are performed on DeepSeekMath-Instruct 7B. The RL training set consists of approximately 144K chain-of-thought-style questions extracted from the GSM8K and MATH subsets of the SFT data; other SFT data is purposefully omitted to examine the effect of RL on evaluation tasks without representation in the RL phase. The reward model training data is built using the methodology described by \citep{wang2023math}. The initial reward model, initialized with DeepSeekMath-Base 7B, is trained with a learning rate of 2e-5. For the GRPO procedure, the policy model utilizes a learning rate of 1e-6, while the KL coefficient is maintained at 0.04. For each question, $64$ candidate answers are sampled. Training employs a maximum sequence length of 1024 tokens and a batch size of 1024. The policy model is updated a single time per exploration cycle. Evaluation of DeepSeekMath-RL 7B is performed on the same suite of benchmarks as DeepSeekMath-Instruct 7B. In this context, GSM8K and MATH (with chain-of-thought reasoning) are considered in-domain tasks, whereas the remaining benchmarks are classified as out-of-domain.

Table \ref{tab:sft_rl_math} presents a comparison of open-source and closed-source models employing both chain-of-thought and tool-based reasoning strategies on English and Chinese evaluation sets. Our analysis reveals: 1) DeepSeekMath-RL 7B achieves accuracy rates of 88.2\% on GSM8K and 51.7\% on MATH when applying chain-of-thought approaches. This level of performance not only exceeds all known open-source models in the 7B to 70B size range but also outperforms most closed-source counterparts. 2) Notably, DeepSeekMath-RL 7B is fine-tuned exclusively on chain-of-thought-style instruction data from GSM8K and MATH, using DeepSeekMath-Instruct 7B as its starting point. Even though its training dataset is limited in scope, DeepSeekMath-RL 7B demonstrates consistent improvements over DeepSeekMath-Instruct 7B on every metric evaluated, highlighting the advantages conferred by reinforcement learning.

\section{Discussion}
In the following section, we outline our principal insights gained from both pre-training and reinforcement learning experiments.

\subsection{Lessons Learnt in Pre-Training}

We begin by summarizing our experiences during the pre-training phase. Unless explicitly mentioned otherwise, all training procedures follow the methodology described in Section \ref{sec:quality-policy}. For reference, "DeepSeekMath Corpus" within this context denotes the 89B-token dataset obtained in the second cycle of our data gathering efforts.

\subsubsection{Code Training Benefits Mathematical Reasoning}

While there is widespread speculation that code training may foster improved reasoning skills, concrete empirical support remains sparse. We seek to contribute some evidence in the realm of mathematics: training with code enhances a model's proficiency in mathematical reasoning, regardless of whether external tools are involved.

To evaluate the effects of code pre-training on mathematical reasoning, we considered both two-stage and one-stage training protocols:

\textbf{Two-Stage Training}

\begin{itemize}[topsep=0pt]
    \item \textbf{Code Training for 400B Tokens $\rightarrow$ Math Training for 150B Tokens}: DeepSeek-LLM 1.3B undergoes training on 400B code tokens, followed by an additional 150B math tokens;
    \item \textbf{General Training for 400B Tokens $\rightarrow$ Math Training for 150B Tokens}: For a comparative study, the model is initially trained with 400B general tokens (sampled from a broad, large-scale corpus curated by DeepSeek-AI) rather than code tokens in the first phase, allowing examination of whether code-centric data provides unique advantages for enhancing mathematical reasoning capabilities.
\end{itemize}

\textbf{One-Stage Training}

\begin{itemize}[topsep=0pt]
    \item \textbf{Math Training for 150B Tokens}: DeepSeek-LLM 1.3B is trained exclusively on 150B math tokens;
    \item \textbf{Training on a mixture of 400B Code Tokens and 150B Math Tokens}: Because code pretraining followed by math finetuning can impair performance on coding tasks, we assess if simultaneous training with a blend of code and math tokens can foster mathematical reasoning abilities while reducing catastrophic forgetting of coding knowledge.
\end{itemize}

\paragraph{Results}
As depicted in Table \ref{tab:code-math} and Table \ref{tab:code-math-general-eval}, downstream results across varied training approaches are provided.

Pretraining on code offers notable gains for program-assisted mathematical reasoning, evident in both two-stage and mixed one-stage settings. Table \ref{tab:code-math} shows that starting with code-focused training in a two-stage process considerably enhances performance on GSM8K and MATH benchmarks utilizing Python-based solutions, while a subsequent math-specific phase leads to even higher gains. Notably, adopting a one-stage mixed curriculum incorporating both code and math data attenuates the catastrophic forgetting seen in sequential two-stage setups and additionally strengthens capabilities in both coding (as seen in Table \ref{tab:code-math-general-eval}) and program-supported mathematical tasks (refer Table \ref{tab:code-math}).

Training with code data also positively influences mathematical problem-solving even without explicit tool usage. Within the two-stage paradigm, preliminary code training delivers measurable improvements and amplifies the efficacy of subsequent math instruction, producing optimal overall performance. On the other hand, fusing code and math in a single training stage appears detrimental for math reasoning tasks that do not leverage tools. A plausible explanation is that the 1.3B parameter DeepSeek-LLM may lack adequate capacity to simultaneously integrate extensive code and math curricula during a unified stage.

\subsubsection{ArXiv Papers Appear Limited in Enhancing Mathematical Reasoning}

Incorporation of arXiv papers as part of mathematical pre-training datasets is a standard approach \citep{minerva,gpt-f,llemma,mathpile}. Nevertheless, there remains a scarcity of systematic investigations into their real effect on models’ mathematical reasoning capabilities. Somewhat unexpectedly, our experimental findings suggest that arXiv papers contribute little to advancing mathematical reasoning abilities. For our evaluation, we examined models across scales, specifically DeepSeek-LLM 1.3B and DeepSeek-Coder-Base-v1.5 7B \citep{deepseek-coder}, and utilized two arXiv-based corpora processed via different methodologies:

\begin{itemize}[topsep=0pt]
    \item \textbf{MathPile} \citep{mathpile}: a curated dataset containing 8.9B tokens, built using various filtering and cleaning heuristics, where scientific arXiv documents comprise more than 85\% of the corpus;
    \item \textbf{ArXiv-RedPajama} \citep{redpajama}: includes the complete set of arXiv LaTeX sources, stripped of preambles, comments, macro definitions, and bibliography sections, amounting to 28.0B tokens.
\end{itemize}

In our experiments, we fine-tuned DeepSeek-LLM 1.3B for 150B tokens and DeepSeek-Coder-Base-v1.5 7B for 40B tokens using each arXiv corpus independently. The results indicate that arXiv-centric data alone does not substantially foster improvements in mathematical reasoning. When these models are exclusively trained on arXiv corpora, performance across diverse mathematical reasoning tasks either stagnates or degrades. Benchmarks assessed range from quantitative reasoning datasets like GSM8K and MATH (Table \ref{tab:arxiv-cot}), multiple-choice evaluations such as MMLU-STEM (Table \ref{tab:arxiv-cot}), to formal mathematics as represented by miniF2F (Table \ref{tab:arxiv_atp}).

Still, these observations should be interpreted with caution due to several caveats. This work does not yet examine:

\begin{itemize}[topsep=0pt]
    \item The possible effect of arXiv-sourced data on particular mathematical tasks outside the present scope, for example, transforming formal mathematical content into informal exposition (theorem informalization);
    \item The impact of blending arXiv data with other types of corpora;
    \item The potential for gains from arXiv papers to emerge at greater model sizes.
\end{itemize}

These aspects remain open and present important directions for subsequent investigations.

\subsection{Reinforcement Learning Insights}
\subsubsection{Toward a Cohesive Framework} Here, we introduce a unified framework that facilitates analysis of different training regimes—including SFT, RFT, DPO, PPO, and GRPO—and empirically investigate influential components within this paradigm. Broadly, the gradient of a training method with respect to parameters $\theta$ can be expressed as follows:

\begin{equation}
    \nabla_{\theta}\mathcal{J}_{\textcolor{red}{\mathcal{A}}}(\theta) = \mathbb{E}[\underbrace{(q,o) \sim \textcolor{red}{\mathcal{D}}}_{Data \ Source}]\left( \frac{1}{|o|} \sum_{t=1}^{|o|}  \underbrace{GC_{{\mathcal{A}}}(q, o, t, \textcolor{red}{\pi_{{rf}}})}_{Gradient \ Coefficient}  \nabla_{\theta}\log \pi_{\theta}(o_t | q, o_{<t})\right).
\label{eq:objective}
\end{equation}

This formulation involves three principal elements: 1) the \textit{Data Source $\mathcal{D}$}, dictating the set of examples for model training; 2) the \textit{Reward Function $\pi_{{rf}}$}, which provides the feedback signal for learning; and 3) the \textit{Algorithm $\mathcal{A}$}, which integrates the data and reward signal into the gradient coefficient $GC$, governing how much the data sample is penalized or rewarded. We consider a range of methods within this shared framework:

\begin{itemize}
    \item  \textbf{Supervised Fine-tuning (SFT)}: This approach involves further training a pretrained model using human-curated supervised data.
    \item \textbf{Rejection Sampling Fine-tuning (RFT)}: Building on SFT, RFT retrains the SFT model on selected outputs—these outputs are generated by the SFT model itself when answering the original SFT questions, and are filtered based on answer accuracy.
    \item \textbf{Direct Preference Optimization (DPO)}: DPO enhances the SFT model by finetuning it using a pairwise loss, where training utilizes expanded sets of outputs sampled from the SFT model.
    \item \textbf{Online Rejection Sampling Fine-tuning (Online RFT)}: Unlike standard RFT, Online RFT commences with the SFT model as the initial policy and iteratively updates it, training with augmented outputs dynamically sampled from the current version of the policy model.
    \item \textbf{PPO/GRPO}: In these methods, the policy model is also initialized with the SFT model, but reinforcement is accomplished by optimizing with outputs continually sampled from the updated policy during training.
\end{itemize}

The elements comprising these approaches are outlined in Table \ref{tab:data_policy}. For a comprehensive exposition of the derivation process, please consult Appendix \ref{app:analysis-rl}.

\paragraph{Observation about Data Source} We categorize the origin of the data into two primary types: online sampling and offline sampling. Online sampling involves deriving training examples directly from the on-policy exploration of the policy being currently trained, whereas offline sampling utilizes data collected via the initial SFT model. Among the methods, RFT and DPO employ an offline approach, whereas Online RFT and GRPO operate with an online sampling paradigm.

According to Figure \ref{fig:rl-analysis}, our analysis reveals that Online RFT achieves notably higher performance than RFT on both evaluation benchmarks. In particular, Online RFT displays similar efficacy to RFT during the initial training phase but establishes a pronounced performance gap in the latter training stages, thereby illustrating the benefits of online data collection. This can be understood as follows: At the onset of training, there is considerable overlap between the actor and the SFT model, resulting in sampled trajectories that are only slightly divergent. As training advances, however, the distinctions in data generated by the evolving actor become more marked, making the real-time acquisition of training data increasingly beneficial.

\paragraph{Observation about Gradient Coefficient} The method utilizes the reward signal to compute the gradient coefficient, which in turn updates the model's parameters. For our experiments, we distinguish between two types of reward functions: `Rule' and `Model'. The `Rule' type assesses response quality through direct evaluation of correctness, while the `Model' type entails learning a reward model that assigns scores to responses, with training examples for the reward model originating from rule-based assessments. A critical distinction highlighted in Equations \ref{eq:GC-RFT} and \ref{eq:GC-GRPO} is that GRPO dynamically modifies its gradient coefficient according to the magnitude of the reward provided by the reward model, thereby enabling variable levels of reinforcement or penalization across responses. Conversely, Online RFT applies uniform reinforcement to all correct responses and does not introduce penalization for incorrect ones, lacking this gradient scaling mechanism.

As illustrated in Figure \ref{fig:rl-analysis}, GRPO outperforms online RFT, underscoring the effectiveness of adjusting the coefficients for positive and negative gradients. Moreover, GRPO+PS achieves better results than GRPO+OS, suggesting that leveraging detailed, step-specific gradient coefficients confers additional advantages. Our investigation extends to iterative RL, where we performed two iterative cycles in the experiments. The results, presented in Figure \ref{fig:iter-rl}, reveal that iterative RL markedly enhances model performance, particularly during the first iteration.

\subsubsection{Why RL Works?} In this study, we apply reinforcement learning using a selected subset from the instruction tuning data, resulting in substantial performance improvements over the instruction-tuned baseline. To clarify the underlying reasons for this efficacy, we assess Pass@K and Maj@K metrics for both the Instruct and RL models across two benchmarks. Figure \ref{fig:maj-pass} demonstrates that RL augments Maj@K accuracy but does not significantly affect Pass@K. This suggests that RL fortifies the model's output distribution, thereby making it more resilient; more precisely, \textbf{the observed improvement appears to arise from elevating the frequency of correct answers among the TopK candidates, rather than fundamental skill gains}. Likewise, \citep{wang2023making} have documented a \textbf{misalignment problem} in the reasoning abilities of SFT models and shown that their performance in reasoning tasks can be boosted via diverse preference alignment methodologies \citep{yuan2023rrhf,song2023preference,wang2023making}.

\subsubsection{How to Achieve More Effective RL?}
\label{sec:effec_rl}
Our experiments demonstrate that RL is highly effective for mathematical reasoning tasks. Additionally, we introduce a unified framework that conceptualizes various representative training strategies, categorizing them as either direct RL approaches or their simplified variants. According to Equation \ref{eq:objective}, this framework is structured around three essential components: Data Source, Algorithm, and Reward Function. We outline some promising directions for future work involving each component.

\paragraph{Data Source} The data source serves as the foundational input for all training procedures. In RL settings, we focus on unlabeled queries coupled with responses generated by the policy model. Our current methodology relies solely on the instruction tuning question set and basic nucleus sampling for output generation. We hypothesize that this might partially account for the RL pipeline's improvement being restricted to Maj@K. Future explorations will include expanding the RL process to out-of-distribution prompts and incorporating \textbf{advanced sampling (decoding) strategies}, such as those leveraging tree-search algorithms \citep{yao2023tree}. Moreover, \textbf{efficient inference techniques} \citep{xia-etal-2023-speculative,leviathan2023fast,kwon2023efficient,xia2024unlocking}—which directly impact the exploration proficiency of policy models—will also be of critical significance.

\paragraph{Algorithms} Algorithms utilize both the input data and the feedback from the reward signal to compute the gradient coefficient used in updating model parameters. As outlined in Equation \ref{eq:objective}, contemporary approaches tend to place considerable \textbf{TRUST} in the reward function's signal when modifying the conditional probability associated with particular tokens. Nevertheless, it is not always feasible to guarantee the accuracy or trustworthiness of the reward signal, especially in highly intricate problem settings. For instance, even PRM800K datasets \citep{lightman2023let}, which are meticulously labeled by skilled annotators, still exhibit around 20\% incorrect annotations\footnote{\url{https://github.com/openai/prm800k/issues/12\#issuecomment-1728491852}}. Consequently, there is a need to examine reinforcement learning algorithms that maintain robustness when exposed to noisy reward feedback. We posit that alignment approaches inspired by the principle of \textbf{WEAK-TO-STRONG} \citep{burns2023weak} will drive significant transformation in learning algorithm design.

\paragraph{Reward Function} The reward function serves as the fundamental origin of the training signal. In reinforcement learning, this typically takes the form of a neural reward model. We identify three major avenues for the advancement of reward models: 1) \textbf{Improving the reward model’s generalization capability.} The reward model should robustly generalize to handle data that falls outside of its training distribution, as well as to complex outputs from advanced decoding strategies; failing this, reinforcement learning risks merely stabilizing rather than genuinely enhancing the LLM's capacity; 2) \textbf{Quantifying the reward model’s uncertainty.} Accurately representing uncertainty can facilitate the interface between an imperfect reward model and weak-to-strong learning approaches; 3) \textbf{Developing efficient, high-fidelity process reward models} that can supply detailed, fine-grained feedback for guiding the reasoning process \citep{lightman2023let,wang2023math}.

\section{Conclusion, Limitation, and Future Work} This work introduces DeepSeekMath, a model that achieves superior results compared to existing open-source systems on the competition-grade MATH benchmark, and approaches the performance level of proprietary counterparts. The training procedure for DeepSeekMath~begins with DeepSeek-Coder-v1.5 7B and includes continual training over 500B tokens, notably incorporating 120B mathematics-specific tokens collected from Common Crawl. Our comprehensive ablation analysis indicates that content from web pages constitutes a rich source of high-quality mathematical data, whereas arXiv contributes less than anticipated. We propose Group Relative Policy Optimization (GRPO), an alternative to Proximal Policy Optimization (PPO), which significantly enhances mathematical reasoning abilities and lowers memory requirements. Experimental findings reveal that GRPO provides performance benefits even when DeepSeekMath-Instruct 7B achieves high benchmark scores. Additionally, we present a unified framework to contextualize various techniques and highlight several promising research directions aimed at advancing reinforcement learning methodologies.

Despite DeepSeekMath’s~strong results on quantitative reasoning tasks, its proficiency in areas such as geometry and theorem proving lags behind leading closed-source models. For example, during preliminary evaluations, the model was unable to address problems involving triangles and ellipses, potentially reflecting biases in the pre-training and fine-tuning data curation. Furthermore, due to its scale, DeepSeekMath’s~performance in few-shot settings trails that of GPT-4. While GPT-4 demonstrates marked improvements with few-shot prompts, DeepSeekMath’s~capabilities remain relatively unchanged between zero-shot and few-shot evaluations. Future directions involve refining our data selection processes to build a higher-quality pre-training dataset, as well as further investigation of the strategies discussed in Section~\ref{sec:effec_rl} to enable more robust reinforcement learning for large language models.